\chapter{Manual de Código}

En este documento se detalla el código fuente completo del sistema de vigilancia distribuido
basado en una red de radares cooperativos. El código se organiza en tres bloques funcionales:
el firmware embebido del nodo radar (C++ sobre ESP32 con FreeRTOS), el servidor central de
orquestación (Python) y las herramientas de visualización y simulación (Python).

Se presupone que el lector tiene conocimientos básicos de programación en C++ y Python,
así como familiaridad con los conceptos de sistemas embebidos y comunicaciones TCP/IP.

% ================================================================
\section{Firmware del Nodo Radar (C++)}
% ================================================================

El firmware se desarrolla en C++ usando el entorno PlatformIO con el núcleo Arduino para
ESP32. La arquitectura multitarea se apoya en FreeRTOS: la tarea \texttt{TaskComms}
gestiona la conexión TCP con el servidor central en el núcleo 0, mientras que
\texttt{TaskRadar} controla el motor paso a paso y el sensor ultrasónico en el núcleo 1.

Para compilar y cargar el firmware basta con abrir la carpeta del proyecto en Visual Studio
Code y pulsar \textbf{Upload} en la barra de PlatformIO, o ejecutar:

\begin{lstlisting}[style=Consola]
pio run --target upload
\end{lstlisting}

% ----------------------------------------------------------------
\subsection{Archivo \texttt{config.h}}
% ----------------------------------------------------------------

Centraliza todas las constantes de configuración del firmware: asignación de pines GPIO,
parámetros del motor (pasos por vuelta, micropasos, ratio de engranaje), rango angular del
barrido, velocidades de homing, credenciales WiFi conocidas y dirección del servidor. Para
adaptar el firmware a hardware distinto basta con modificar únicamente este archivo.

\lstinputlisting[
  style=C++-color,
  caption={config.h — Constantes de configuración globales del firmware.},
  label={lst:config}
]{src/config.h}

% ----------------------------------------------------------------
\subsection{Archivo \texttt{comms\_protocol.h}}
% ----------------------------------------------------------------

Define el protocolo binario de comunicación entre los nodos y el servidor central. Cada
mensaje se compone de una cabecera de tres bytes (\texttt{type\,+\,length}) seguida de
una carga útil de longitud fija. Las estructuras se declaran con \texttt{\#pragma~pack(1)}
para eliminar el relleno del compilador y garantizar que la representación en memoria sea
idéntica en el ESP32 y en el servidor Python.

\lstinputlisting[
  style=C++-color,
  caption={comms\_protocol.h — Opcodes y estructuras del protocolo binario.},
  label={lst:protocol}
]{src/comms_protocol.h}

% ----------------------------------------------------------------
\subsection{Archivo \texttt{system.h}}
% ----------------------------------------------------------------

Declara la máquina de estados (\texttt{SystemState}) y el gestor de sistema
(\texttt{SystemManager}), implementado como \textit{singleton} de FreeRTOS. El
\texttt{SystemManager} mantiene el estado global compartido entre ambas tareas: las colas
de mensajes, el identificador del nodo, el ángulo lógico actual y los perfiles de hardware
específicos de cada unidad (dirección del motor, nivel de activación del reed switch).

\lstinputlisting[
  style=C++-color,
  caption={system.h — FSM y gestor de estado global compartido entre tareas.},
  label={lst:system_h}
]{src/system.h}

% ----------------------------------------------------------------
\subsection{Archivo \texttt{system.cpp}}
% ----------------------------------------------------------------

Implementa los métodos del \texttt{SystemManager}: la inicialización de las colas FreeRTOS
y el método de transición de estado, que imprime un mensaje de diagnóstico por el puerto
serie en cada cambio.

\lstinputlisting[
  style=C++-color,
  caption={system.cpp — Inicialización de colas FreeRTOS y transiciones de estado.},
  label={lst:system_cpp}
]{src/system.cpp}

% ----------------------------------------------------------------
\subsection{Archivo \texttt{task\_radar.h}}
% ----------------------------------------------------------------

Define los tipos de mensajes intercambiados entre \texttt{TaskComms} y \texttt{TaskRadar}
a través de las colas FreeRTOS: \texttt{HwCommand} (órdenes enviadas al hardware) y
\texttt{HwResult} (resultados devueltos desde el hardware). Incluye la declaración de la
clase \texttt{RadarHardware}, que encapsula el control del motor paso a paso y del sensor
ultrasónico.

\lstinputlisting[
  style=C++-color,
  caption={task\_radar.h — Interfaz de la tarea hardware y de RadarHardware.},
  label={lst:radar_h}
]{src/task_radar.h}

% ----------------------------------------------------------------
\subsection{Archivo \texttt{task\_comms.h}}
% ----------------------------------------------------------------

Cabecera mínima de la tarea de comunicaciones. Declara únicamente el prototipo de
\texttt{TaskComms}, que es el punto de entrada de la tarea FreeRTOS responsable del
ciclo TCP con el servidor.

\lstinputlisting[
  style=C++-color,
  caption={task\_comms.h — Prototipo de la tarea de comunicaciones.},
  label={lst:comms_h}
]{src/task_comms.h}

% ----------------------------------------------------------------
\subsection{Archivo \texttt{main.cpp}}
% ----------------------------------------------------------------

Contiene \texttt{setup()} y \texttt{loop()} del núcleo Arduino. \texttt{setup()} inicializa
el hardware, aplica el perfil del nodo según su dirección MAC (tabla \texttt{PROFILES[]}),
registra el manejador de eventos WiFi y lanza el proceso de conexión.

\texttt{loop()} implementa la FSM de alto nivel: parpadeo del LED de estado según la
conectividad, reintento de conexión WiFi con escaneo periódico y creación de las tareas
FreeRTOS en el núcleo adecuado cuando el estado del sistema lo requiere.

El bloque condicional \texttt{\#if~DEBUG\_HARDWARE\_TEST~==~1} activa un modo de
diagnóstico interactivo que verifica el sensor ultrasónico, el reed switch y el motor paso
a paso de forma autónoma, sin necesidad de servidor.

\lstinputlisting[
  style=C++-color,
  caption={main.cpp — Inicialización, perfiles de nodo y bucle principal de la FSM.},
  label={lst:main}
]{src/main.cpp}

% ----------------------------------------------------------------
\subsection{Archivo \texttt{task\_radar.cpp}}
% ----------------------------------------------------------------

Implementa la tarea \texttt{TaskRadar} y los métodos de \texttt{RadarHardware}. La tarea
espera bloqueada en la cola de comandos y despacha cada \texttt{HwCommand}: \texttt{HW\_CMD\_HOME}
invoca \texttt{goHome()}, que mueve el motor al ángulo cero y verifica la posición con el reed
switch; \texttt{HW\_CMD\_MOVE} y \texttt{HW\_CMD\_EXECUTE\_SLOT} llaman a \texttt{moveToAngle()},
que calcula y ejecuta los pasos necesarios; \texttt{HW\_CMD\_SYNC\_POS} actualiza la
variable de ángulo sin mover el motor (útil tras una reconexión). El resultado se envía a
\texttt{TaskComms} mediante la cola de resultados.

\lstinputlisting[
  style=C++-color,
  caption={task\_radar.cpp — Control del motor paso a paso, homing y medición ultrasónica.},
  label={lst:radar_cpp}
]{src/task_radar.cpp}

% ----------------------------------------------------------------
\subsection{Archivo \texttt{task\_comms.cpp}}
% ----------------------------------------------------------------

Implementa la tarea \texttt{TaskComms}, que gestiona el ciclo de vida completo de la
conexión TCP: establecimiento de la conexión, \textit{handshake} HELLO, recepción y
procesamiento de todos los tipos de mensaje del protocolo (inicio de supertrama, asignación
de ranura, solicitud de reporte y petición de ángulo) y envío de las respuestas. Incluye
los mecanismos de recuperación ante fallos: \textit{timeout} de servidor inactivo de 5\,s,
contador de fallos TCP y reinicio del ESP32 al superar el límite de fallos consecutivos.

\lstinputlisting[
  style=C++-color,
  caption={task\_comms.cpp — Gestión del protocolo TCP y máquina de estados de comunicación.},
  label={lst:comms_cpp}
]{src/task_comms.cpp}

% ================================================================
\section{Servidor Central (Python)}
% ================================================================

El servidor se desarrolla en Python 3 y se ejecuta en un ordenador de propósito general
conectado a la misma red WiFi que los nodos. Implementa el bucle de supertrama TDMA, el
algoritmo de asignación dinámica de ranuras con reutilización espacial, la lógica de
seguimiento cooperativo, la trilateración y la detección de desconexiones. La comunicación
con los nodos se realiza mediante sockets TCP; la comunicación con el visualizador, mediante
datagramas UDP. No requiere dependencias externas más allá de la biblioteca estándar de Python.

Para arrancar el servidor de forma autónoma (sin el visualizador):

\begin{lstlisting}[style=Consola]
cd server/
python server_central.py
\end{lstlisting}

% ----------------------------------------------------------------
\subsection{Archivo \texttt{server\_central.py}}
% ----------------------------------------------------------------

Archivo principal del servidor. Contiene la clase \texttt{RadarServer}, que gestiona el
registro de nodos mediante mDNS con el \textit{hostname} \texttt{radar-server}, el cálculo
de pares conflictivos (zona sucia), la ejecución del ciclo de supertrama TDMA y el algoritmo
de trilateración cooperativa. El estado angular de todos los nodos se persiste en disco en
\texttt{config/radar\_state.json} para que la sesión pueda reanudarse.

\lstinputlisting[
  style=Python-color,
  caption={server\_central.py — Servidor central de orquestación TDMA y trilateración.},
  label={lst:server}
]{server/server_central.py}

% ================================================================
\section{Herramientas e Interfaz Gráfica (Python)}
% ================================================================

Las herramientas se desarrollan en Python 3 usando exclusivamente bibliotecas incluidas en
la distribución estándar (\texttt{tkinter}, \texttt{socket}, \texttt{threading},
\texttt{subprocess}, \texttt{json}). Las únicas dependencias externas son \texttt{zeroconf}
e \texttt{ifaddr}, necesarias para el anuncio y descubrimiento mDNS:

\begin{lstlisting}[style=Consola]
pip install zeroconf ifaddr
\end{lstlisting}

% ----------------------------------------------------------------
\subsection{Archivo \texttt{radar\_visualizer.py}}
% ----------------------------------------------------------------

Interfaz gráfica principal del sistema. Implementa dos modos de operación: el \textbf{Modo
Diseño}, que permite configurar la geometría del despliegue y los umbrales DEFCON antes
de iniciar la sesión, y el \textbf{Modo Telemetría}, que representa en tiempo real las
detecciones recibidas del servidor mediante datagramas UDP. Gestiona también el arranque y
la parada del servidor central como subproceso y la persistencia de sesión entre ejecuciones.

Para ejecutarlo:

\begin{lstlisting}[style=Consola]
cd tools/
python radar_visualizer.py
\end{lstlisting}

\lstinputlisting[
  style=Python-color,
  caption={radar\_visualizer.py — Interfaz gráfica de visualización y control del sistema.},
  label={lst:visualizer}
]{tools/radar_visualizer.py}

